\par\Needspace{6\baselineskip}\medskip
\hypertarget{native-209-article155}{}\pdfbookmark[2]{155.}{article-155}
\begin{center}155.\end{center}
Quo clarius elucescat, quanta praecisione elementa inventa gaudeant, locum medium ex ipsis computabimus. Pro Oct.\ $17{,}415\,011$ anomalia media invenitur $=332^{0}\,28'\,54''{,}76$, hinc vera $315^{0}\,1'\,23''{,}02$ atque $\log r'=0{,}325\,9877$ (vid.\ \SourceLeaf{210}exempla artt.\ 13, 14); illa aequalis esse deberet anomaliae verae in observatione prima auctae angulo $2f''$, vel anomaliae verae in observatione tertia diminutae angulo $2f$, i.\ e.\ $=315^{0}\,1'\,22''{,}98$; logarithmus radii vectoris vero $=0{,}325\,9878$: differentiae pro nihilo habendae sunt. Si calculus pro observatione media usque ad locum geocentricum continuatur, numeri resultant ab observatione paucis tantum minuti secundi partibus centesimis deviantes (art.\ 63), quales differentiae ab erroribus inevitabilibus e tabularum praecisione limitata oriundis quasi absorbentur.

Exemplum praecedens summa praecisione ideo tractavimus, ut appareat, quam facile per methodum nostram solutio quam accuratissima obtineri possit. In ipsa praxi raro opus erit, hunc typum aeque anxie imitari: plerumque sufficiet, \textit{sex} figuras decimales ubique adhibere, et in exemplo nostro secunda iam hypothesis praecisionem haud minorem, primaque praecisionem abunde sufficientem suppeditavisset. Haud ingratam fore lectoribus censemus comparationem elementorum ex hypothesi tertia erutorum cum iis, quae prodeunt, si hypothesis secunda vel adeo prima perinde ad eundem scopum adhibitae fuissent. Haec tria elementorum systemata in schemate sequente exhibemus:
\begin{center}
\begingroup\footnotesize\setlength{\tabcolsep}{2pt}\renewcommand{\arraystretch}{1.2}
\begin{tabular}{|p{1.35in}|r|r|r|}\hline
& ex hypothesi III & ex hypothesi II & ex hypothesi I\\\hline
Epocha longit.\ med.\ 1805 & $41^{0}\,52'\,21''{,}68$ & $41^{0}\,52'\,18''{,}40$ & $42^{0}\,12'\,37''{,}83$\\
motus medius diurnus & $824''{,}7990$ & $824''{,}7978$ & $823''{,}5026$\\
perihelium & $52^{0}\,18'\,9''{,}30$ & $52^{0}\,18'\,6''{,}66$ & $52^{0}\,41'\,9''{,}81$\\
$\varphi$ & $14\quad12\quad1{,}87$ & $14\quad11\quad59{,}94$ & $14\quad24\quad27{,}49$\\
logar.\ semiaxis maioris & $0{,}422\,4389$ & $0{,}422\,4394$ & $0{,}422\,8944$\\
nodus ascendens & $171^{0}\,7'\,48''{,}73$ & $171^{0}\,7'\,49''{,}15$ & $171^{0}\,5'\,48''{,}86$\\
inclinatio orbitae & $13\quad6\quad44{,}10$ & $13\quad6\quad45{,}12$ & $13\quad2\quad37{,}50$\\\hline
\end{tabular}\endgroup
\end{center}

Computando locum heliocentricum in orbita pro observatione media per secundum elementorum systema, invenitur error logarithmi radii vectoris $=0$, error longitudinis in orbita $=0''{,}03$; computando vero istum locum per systema ex hypothesi prima derivatum prodit error logarithmi radii vectoris $=0{,}000\,0002$, error longitudinis in orbita $=1''{,}31$. Continuando vero calculum usque ad locum geocentricum invenitur
\SourceLeaf{211}
\begin{center}
\begingroup\small\setlength{\tabcolsep}{4pt}\renewcommand{\arraystretch}{1.2}
\begin{tabular}{|l|r|r|}\hline
& ex hypothesi II & ex hypothesi I\\\hline
Longitudo geocentrica & $352^{0}\,34'\,22''{,}26$ & $352^{0}\,34'\,19''{,}97$\\
error & $0{,}14$ & $2{,}15$\\
latitudo geocentrica & $-6\quad21\quad55{,}06$ & $-6\quad21\quad54{,}47$\\
error & $0{,}01$ & $0{,}60$\\\hline
\end{tabular}\endgroup
\end{center}
